%% thread.tex
%%
%% Copyright (C) 2007-2018 Simone Piccardi.  Permission is granted to
%% copy, distribute and/or modify this document under the terms of the GNU Free
%% Documentation License, Version 1.1 or any later version published by the
%% Free Software Foundation; with the Invariant Sections being "Un preambolo",
%% with no Front-Cover Texts, and with no Back-Cover Texts.  A copy of the
%% license is included in the section entitled "GNU Free Documentation
%% License".
%%

\chapter{I thread}
\label{cha:threads}


\itindbeg{thread} 

Tratteremo in questo capitolo un modello di programmazione multitasking,
quello dei \textit{thread}, alternativo al modello classico dei processi,
tipico di Unix. Ne esamineremo le caratteristiche, vantaggi e svantaggi, e le
diverse realizzazioni che sono disponibili per Linux; nella seconda parte
tratteremo in dettaglio quella che è l'implementazione principale, che fa
riferimento all'interfaccia standardizzata da POSIX.1e.


\section{Introduzione ai \textit{thread}}
\label{sec:thread_intro}

Questa prima sezione costituisce una introduzione ai \textit{thread} e
tratterà i concetti principali del relativo modello di programmazione,
esamineremo anche quali modelli sono disponibili per Linux, dando una breve
panoramica sulle implementazioni alternative.


\subsection{Una panoramica}
\label{sec:thread_overview}

% riferimenti
% http://vergil.chemistry.gatech.edu/resources/programming/threads.html
% http://math.arizona.edu/~swig/documentation/pthreads/
% http://www.humanfactor.com/pthreads/

Il modello classico dell'esecuzione dei programmi nei sistemi Unix, illustrato
in sez.~\ref{cha:process_interface}, è fondato sui processi. Il modello nasce
per assicurare la massima stabilità al sistema e prevede una rigida
separazione fra i diversi processi, in modo che questi non possano disturbarsi
a vicenda. 

Le applicazioni moderne però sono altamente concorrenti, e necessitano quindi
di un gran numero di processi; questo ha portato a scontrarsi con alcuni
limiti dell'architettura precedente. In genere i fautori del modello di
programmazione a \texttt{thread} sottolineano due problemi connessi all'uso
dei processi:
\begin{itemize}
\item
\item 
\end{itemize}



\subsection{\textit{Thread} e processi}
\label{sec:thread_process}

Per un utilizzo effettivo dei \textit{thread} è sempre opportuno capire se
questi sono davvero adatti allo scopo che ci si pone.


\subsection{Implementazioni alternative}
\label{sec:thread_other}

Vedremo nella prossima sezione le caratteristiche del supporto per i
\textit{thread} fornita dal kernel, ma esistono diversi possibili approcci
alle modalità in cui questi possono essere realizzati. 

% TODO cenni su pth e le implementazioni in userspace


\section{I \textit{thread} e Linux}
\label{sec:linux_thread}

In questa sezione tratteremo le implementazioni dei \textit{thread}
disponibili con Linux che ha visto un radicale cambiamento nel passaggio dalla
serie 2.4 alla serie 2.6, che ha portato alla versione attuale. 

\subsection{I \textit{LinuxThread}}
\label{sec:linux_old_thread}


\subsection{La \textit{Native Thread Posix Library}}
\label{sec:linux_ntpl}





% http://www.gnu.org/software/pth/


\section{Posix \textit{thread}}
\label{sec:thread_posix_intro}


Tratteremo in questa sezione l'interfaccia di programmazione con i
\textit{thread} standardizzata dallo standard POSIX 1.c, che è quella che è
stata seguita anche dalle varie implementazioni dei \textit{thread} realizzate
su Linux, ed in particolare dalla \textit{Native Thread Posix Library} che è
stata integrata con i kernel della serie 2.6 e che fa parte a pieno titolo
della \acr{glibc}.


\subsection{Una panoramica}
\label{sec:pthread_overview}


\subsection{La gestione dei \textit{thread}}
\label{sec:pthread_management}


Benché la funzione sia utilizzabile anche con i processi, tanto che a partire
dalla versione 2.3 della \acr{glibc} viene a sostituire \func{\_exit} (tramite
un \textit{wrapper} che la utilizza al suo posto) per la terminazione di tutti
i \textit{thread} di un processo si deve usare la funzione di sistema
\func{exit\_group}, il cui prototipo è:

\begin{funcproto}{
\fhead{linux/unistd.h}
\fdecl{void exit\_group(int status)}
\fdesc{Termina tutti i \textit{thread} di un processo.} 
}
{La funzione non ha errori e pertanto non ritorna.}
\end{funcproto}

La funzione è sostanzialmente identica alla \textit{system call} \func{\_exit}
ma a differenza di quest'ultima, che termina solo il \textit{thread}
chiamante, termina tutti \textit{thread} del processo. 


\section{La sincronizzazione dei \textit{thread}}
\label{sec:pthread_sync}

\subsection{I \textit{mutex}}
\label{sec:pthread_mutex}


\subsection{Le variabili di condizione}
\label{sec:pthread_cond}


\itindend{thread} 


\subsection{I \textit{thread} e i segnali.}
\label{sec:thread_signal}

% TODO trattare tkill e tgkill per l'invio di segnali a thread, fare un
% capitolo apposito su thread e segnali




% TODO troppe cose, ma segue list di notizie correlate
% aggiunta rt_tgsigqueueinfo con il kernel 2.6.31



% LocalWords:  thread multitasking POSIX sez Posix Library kernel glibc mutex


%%% Local Variables: 
%%% mode: latex
%%% TeX-master: "gapil"
%%% End: 

